\section{The first coupled vector correction}\label{sec:coupled}

\subsection{Action convention and the corrected background}
For this calculation use unit-scale internal coordinates and remove the
common gravitational prefactor. Set $\eta=\alpha'/L^2$ and use positive
forms $\BF=-\tfrac12\Tr_{32}$ and $\BR=-\tfrac12\Tr_T$ on
the compact gauge and internal Euclidean rotation algebras. The
Lorentzian extension of the tangent trace is understood in the
ten-dimensional action. Our convention is
\begin{equation}\label{eq:heterotic-action}
 \begin{split}
 S=\int e^{-2\Phi}\sqrt{-G}\Bigl[&
 R+4|d\Phi|^2-\frac1{12}H_{MNP}H^{MNP}\\
 &-\frac\eta4\bigl(\BF(F_{MN},F^{MN})-\BR(R_{MN},R^{MN})\bigr)
 \Bigr].
 \end{split}
\end{equation}
Use the field choice in which no additional independent purely bosonic
operator at order $\alpha'^2$ is required beyond the recursively defined
torsion and composite-connection terms \cite{bdr,aqs}. This statement
does not discard their order-$\alpha'^2$ contributions. The tangent
connection is the composite Hull connection
$\Theta_H=\Theta_{\rm LC}+\tfrac12H$, where
$H^a{}_b=H^a{}_{bM}dx^M$ in an orthonormal frame, in the convention
of \cite[section~2.1]{mcorist_yin}. It reduces to Levi--Civita at leading order.
Fix Chern--Simons signs by
\begin{align}
 dH&=\frac\eta2\bigl(\BR(R\wedge R)-\BF(F\wedge F)\bigr),
 \label{eq:Bianchi}\\
 \delta H&=d\mathcal B-\eta\BF(\delta\mathcal A\wedge F)
                  +\eta\BR(\delta\Theta\wedge R).
 \label{eq:Hvariation}
\end{align}
$\mathcal B$ denotes the gauge-covariantly shifted two-form variation.
These choices fix relative normalizations rather than combining formulas
with incompatible trace or dilaton conventions.

For $t=(a,b)$ let $F_t$ be the full leading ASD curvature. Through first
order the smooth background is
\begin{equation}\label{eq:background}
 G_S=(1+\eta f_t)g_0,\qquad
 \Phi=\Phi_0+\tfrac12\eta f_t,\qquad H=-\eta *_0df_t,
\end{equation}
where $f_t$ has zero mean and solves
\begin{equation}\label{eq:Poisson}
 \nabla^m\nabla_m f_t
 =\frac14\bigl(\BR(R_{mn},R^{mn})-\BF(F_{t,mn},F_t^{mn})\bigr).
\end{equation}
Here $\nabla^2$ has the divergence sign; the nonnegative scalar
Laplacian is $\Delta_0=-\nabla^2$. The source has mean zero by
\eqref{eq:instanton}. Scalar elliptic inversion gives a unique
solution, with uniform spatial bounds on a sufficiently small family
of reconstructed connections. These connections have uniformly bounded
spatial coefficients in fixed gauge even though the small spectral
ratios degenerate. Parameter derivative bounds used later are only on
fixed regular patches. For sufficiently small $\eta$, $1+\eta f_t>0$.

The normalization can be checked directly. In four dimensions the
linear conformal Ricci tensor is
$-\nabla_m\nabla_nf-\tfrac12g_{mn}\nabla^2f$; the dilaton
Hessian adds $\nabla_m\nabla_nf$. For an ASD curvature,
$\BF(F_{mp},F_n{}^p)=\tfrac14g_{mn}\BF(F_{pq},F^{pq})$.
The Einstein equation from \eqref{eq:heterotic-action} gives
\eqref{eq:Poisson}. Also
$\BF(F\wedge F)=-\tfrac12\BF(F_{mn},F^{mn})dV$ and
$d(-*df)=-\nabla^2 f\,dV$, so the Bianchi identity gives the same
equation. The dilaton equation gives the same trace condition, and
the two-form equation holds because the leading weighted dual flux
is exact. The gauge equation
$D_{\mathcal A}(e^{-2\Phi}*_4F)-e^{-2\Phi}F\wedge *_4H=0$
also cancels through this order, using $*_4F=-F$ and
$*_4H=\eta df$. The corresponding conformal/dilaton ansatz is the usual
first-order supersymmetric K3 construction \cite{gpt}; ASD persists
under conformal rescaling. It is a solution through the computed
order, not an existence assertion for an exact finite-$\eta$ background.

\subsection{The two heavy spin-one sectors}
Let $P_t$ be the whole 96-dimensional low real spectral projection
of $d_{\mathcal A_t}^*d_{\mathcal A_t}$ on the full orthogonal
adjoint. Put $A_t=P_td_{\mathcal A_t}^*d_{\mathcal A_t}P_t$.
This space includes all persistent massless vectors. Define
\begin{equation}\label{eq:TJ}
 T_tp=\BF(p,F_t)\in\Omega^2(S),\qquad
 (J_tp)_m=\BF(F_{t,mn},D^np),\qquad
 I_tv=\iota_vF_t.
\end{equation}
The Yang--Mills equation and invariance of the trace give
$J_t=d^*T_t=I_t^*d_{\mathcal A_t}$, and hence $d^*J_t=0$.
In components the latter follows by differentiating
$\BF(F_{mn},D^np)$; the curvature commutator contracts to zero.

Set $s=\sqrt\eta$ and canonically normalize the gauge fluctuation by
$p=s\,\delta\mathcal A_\mu$. Let $v$ be the metric vector in the
coframe $dy^m+v^m_\mu dx^\mu$, and let $\beta$ be the shifted
two-form vector. After the diffeomorphism and form constraints their
profiles are co-closed real one-forms. At leading order both have
ordinary $L^2$ kinetic norm and squared-mass operator
$\Delta_1=d^*d$ on co-closed one-forms. The external field-strength
coefficient is $-1/4$ for each. For the two-form, the three placements
of the internal index in $H_{\mu\nu m}$ turn its contribution to
$-H^2/12$ into $-H_{\mu\nu m}H^{\mu\nu m}/4$, giving the
external vector kinetic term. The three placements of the external
index in $H_{\mu mn}$ instead give its mass term; the factor
$1/2$ in the internal two-form norm yields the usual Proca
normalization. For the metric the kinetic coefficient follows from the
quadratic Kaluza--Klein curvature formula; on a Ricci-flat base its
Killing-operator form reduces, on co-closed forms, to $\|dv\|^2$.
K3 has no harmonic one-forms, and a Killing field on a compact
Ricci-flat K3 would be parallel. Thus both sectors have a fixed
positive gap.

With $I_tv=\iota_vF_t$ and the plus coframe just specified, the
horizontal mixed gauge curvature is
$-d_{\mathcal A}\delta\mathcal A_\mu-I_tv_\mu$.
The terms determining the first low-vector coefficient are therefore
\begin{equation}\label{eq:mixing-squares}
 Q=\|d_{\mathcal A}p+sI_tv\|^2+\|dv\|^2
                      +\|d\beta-sT_tp\|^2,
 \qquad K=\|p\|^2+\|v\|^2+\|\beta\|^2
\end{equation}
at leading heavy order. Reversing the coframe convention reverses
the metric mixing sign and leaves its Schur term unchanged. Exact
one-form gauge directions are not additional transverse vector
species. Harmonic two-form components of $T_tp$, in contrast, must
remain in the direct norm in \eqref{eq:mixing-squares}.

The conformal background contributes gauge kinetic weight
$(1+\eta f_t)$: its four-dimensional volume factor and dilaton
factor leave one power of $1+\eta f_t$. The extra inverse internal
metric in the gauge mass term cancels that power to first order.
Thus the direct gauge correction to the pencil is
$T_t^*T_t-\lambda f_t$.

Eliminating both complete leading heavy towers gives
\begin{equation}\label{eq:effectivecoefficient}
 E_t(\lambda)=P_t\left[T_t^*T_t-\lambda f_t
           -2J_t^*(\Delta_1-\lambda)^{-1}J_t\right]P_t.
\end{equation}
Each factor of two in this expression comes from one of the two
distinct, equally normalized heavy sectors. No finite KK cutoff is
taken. Leading elliptic estimates give
$\|p\|_{H^k}\le C_k\|p\|$ for $p\in\operatorname{im}P_t$,
uniformly in $t$. Together with the one-form gap and the uniform
Poisson bound, this proves, on a fixed interval below half that gap,
\begin{equation}\label{eq:Ebound}
 \|E_t(\lambda)\|+\|\partial_\lambda E_t(\lambda)\|\le C.
\end{equation}
There is no inverse small gauge eigenvalue in this bound.

For clarity, the order count also covers the terms not visible in
\eqref{eq:mixing-squares}. Metric--two-form heavy mixing from the
background $H$ and first-order changes of their operators are of
order $\eta$. Between two gauge--heavy mixings of order
$\sqrt\eta$ they first affect the low pencil at order $\eta^2$.
The $\eta\|I_tv\|^2$ term has the same property. A gauge-vector
induced torsion variation is of order $\sqrt\eta$ in canonical
variables, so its quadratic contribution through
$\eta R(\Theta)^2$ starts at order $\eta^2$.
Quadratic gauge Chern--Simons terms have at least two external legs
and cannot contract with the purely internal background to supply
an earlier diagonal term. Lorentz Chern--Simons corrections to a
gauge--heavy mixing start at order $\eta^{3/2}$. Finally the
off-diagonal coupling to the gapped gauge complement starts at
order $\eta$, and its Schur contribution at order $\eta^2$.
The absence of a further independent bosonic $\alpha'^2$ operator
in the chosen scheme matters here: canonical gauge normalization
removes one power of $\eta$ from a quadratic gauge term.
We extract a perturbative coefficient; we do not treat the raw
curvature-squared action as a positive bounded perturbation of an
infinite-dimensional second-order form.

\subsection{Hodge cancellation, positivity and the error bound}
Let $\HH_2$ denote harmonic projection on real two-forms. For every
adjoint zero-form $p$, ASD gives $*T_tp=-T_tp$. Decompose
$T_tp=h+e+c$ into harmonic, exact and coexact parts. Hodge star
preserves $h$ and exchanges $e,c$ isometrically, so
$\|e\|=\|c\|$. Moreover
$e=d\Delta_1^{-1}d^*T_tp$. Hence
\begin{equation}\label{eq:Hodgeidentity}
 \|T_tp\|^2-2\langle J_tp,\Delta_1^{-1}J_tp\rangle
                   =\|\HH_2T_tp\|^2.
\end{equation}
This elementary identity is a useful check on the two heavy
normalizations. Polarization and the resolvent identity rewrite
\eqref{eq:effectivecoefficient} as
\begin{align}\label{eq:Epositive}
 E_t(\lambda)&=P_tT_t^*\HH_2T_tP_t-\lambda Z_t(\lambda),\\
 Z_t(\lambda)&=P_t\left[f_t+
 2J_t^*\Delta_1^{-1}(\Delta_1-\lambda)^{-1}J_t\right]P_t.
 \nonumber
\end{align}

Define the first-order effective model by
\begin{equation}\label{eq:pencil}
 F_t^{[1]}(\lambda,\eta)=A_t-\lambda I+\eta E_t(\lambda),
 \qquad m^{2,[1]}=L^{-2}\lambda.
\end{equation}
For small $\eta$, its zero-energy mass form is nonnegative,
and its kinetic form $I+\eta Z_t(\lambda)$ is positive on the
fixed low interval. Also $\partial_\lambda F_t^{[1]}=-I+O(\eta)$
is negative definite there. These assertions concern a finite
96-dimensional pencil with an exactly specified coefficient.

Every parallel generator of the unbroken algebra has
$d_{\mathcal A_t}p=0$ and $T_tp=0$ pointwise. In the central
unitary directions the latter follows from determinant-flat
curvature on each stable factor; it also holds in the trivial
$\SO(14)$ block. On an axis the separate rank-three determinants
remain trivial because $\operatorname{Pic}^0(S)=0$.
Consequently
\begin{equation}\label{eq:kernel}
 \ker F_t^{[1]}(0,\eta)=\ker A_t.
\end{equation}
There are exactly 93, 94 or 96 zero roots on the corresponding
strata. Positivity precludes an extra kernel outside $\ker A_t$.

\begin{proposition}\label{prop:coupled}
For sufficiently small $a,b,\eta\ge0$, the pencil
\eqref{eq:pencil} has 96 real low roots, counted with
multiplicity, with kernels as in \eqref{eq:kernel}. Its three
positive generic branches satisfy \eqref{eq:intro-corrected}.
\end{proposition}
\begin{proof}
Choose a fixed positive cutoff below the one-form and gauge
complementary gaps and then shrink the parameter neighborhood so
that the whole leading low cluster lies well below it. The
pencil is positive at a small fixed negative endpoint and
negative at the positive endpoint, and decreases strictly
between them. Its ordered eigenvalue functions therefore have
exactly one zero each. Equation~\eqref{eq:Ebound} and the
min--max inequalities displace each root from the corresponding
eigenvalue of $A_t$ by at most $C\eta$. Divide by the positive
finite eigenvalue, whose minimum is a fixed constant times
$b^2$, and apply Proposition~\ref{prop:vectors}. Exact zero
branches are treated by \eqref{eq:kernel}, not by division.
\end{proof}

The formal full expression is
$F=A_t-\lambda I+\eta E_t(\lambda)+\Oformal(\eta^2)$.
Here $\Oformal$ records uncomputed higher coefficients; it is
not a uniform numerical bound on an exact finite-$\eta$ pole.
Solving \eqref{eq:pencil} retains its energy dependence and
thereby resums some higher powers by convention. Those powers
are not additional calculated string coefficients. Along
$b=\varepsilon^p$, $a=\varepsilon$, $p>1$, the sufficient
condition $\eta=o(\varepsilon^{2p+2})$ refers precisely to
the model estimate just proved.
